 %-------------------------------------- COMIENZA descripción del caso de uso.

  \begin{UseCase}{CU05}{Agregar Vacuna}{
      Permite al Propietario registrar una nueva vacuna en el historial médico de una de sus mascotas, especificando el tipo de vacuna, fechas de aplicación
  y vigencia.
  }
      \UCitem{Versión}{\color{Gray}1.0}
      \UCitem{Autor}{\color{Gray}Jared}
      \UCitem{Supervisa}{\color{Gray}Ulises velez saldaña}
      \UCitem{Actor}{\hyperlink{Propietario}{Propietario}}
      \UCitem{Propósito}{Mantener un registro completo y actualizado de las vacunaciones de una mascota para asegurar su salud, cumplir con los requisitos
  del hotel y llevar un control de las fechas de revacunación.}
      \UCitem{Entradas}{
          \begin{itemize}
              \item ID de la mascota (desde la URL).
              \item ID de la vacuna (opcional, del catálogo).
              \item Nombre de la vacuna (opcional, si no se selecciona del catálogo).
              \item Fecha de aplicación.
              \item Fecha de vigencia.
              \item Nombre del veterinario.
              \item Notas adicionales.
          \end{itemize}
      }
      \UCitem{Origen}{Teclado / Mouse}
      \UCitem{Salidas}{
          \begin{itemize}
              \item Mensaje de confirmación del registro.
              \item Redirección a la pantalla de detalle de la mascota con la lista de vacunas actualizada.
          \end{itemize}
      }
      \UCitem{Destino}{Pantalla y Base de Datos}
      \UCitem{Precondiciones}{
          \begin{itemize}
              \item El Propietario ha iniciado sesión en el sistema.
              \item El Propietario se encuentra en la \IUref{IU-PetDetail}{Pantalla de Detalle de Mascota}.
          \end{itemize}
      }
      \UCitem{Postcondiciones}{Se crea un nuevo registro en la tabla vacunas mascotas de la base de datos, asociado a la mascota correspondiente.}
      \UCitem{Errores}{
          \begin{itemize}
              \item \textbf{E1}: Datos del formulario inválidos (ej. formato de fecha incorrecto).
              \item \textbf{E2}: El usuario no está autenticado o la sesión ha expirado.
              \item \textbf{E3}: La mascota especificada no existe o no pertenece al usuario.
              \item \textbf{E4}: Error inesperado del servidor al intentar guardar el registro.
          \end{itemize}
      }
      \UCitem{Tipo}{Secundario}
  \end{UseCase}
  %--------------------------------------
  % Trayectoria Principal
  %--------------------------------------
  \begin{UCtrayectoria}{Principal}
      \UCpaso[\UCactor] Desde la \IUref{IU-PetDetail}{Pantalla de Detalle de Mascota}, oprime el botón \IUbutton{Agregar Vacuna}.
      \UCpaso El sistema despliega la \IUref{IU-AddVaccineForm}{Formulario de Registro de Vacuna}.
      \UCpaso[\UCactor] Ingresa los datos de la vacunación en los campos correspondientes. \label{CU05-FillForm}
      \UCpaso[\UCactor] Oprime el botón \IUbutton{Guardar Vacuna}. \label{CU05-Submit}
      \UCpaso El sistema valida que los datos ingresados tengan el formato correcto con base en la regla \BRref{BR-VaccineValidation}{Validación de Datos de Vacunación}. \Trayref{A}.
      \UCpaso El sistema envía una petición POST /api/pet-vaccinations/:mascotaId al servidor.
      \UCpaso El sistema verifica que la sesión del usuario sea válida con base en la regla \BRref{BR-AuthCheck}{Verificación de Autenticación}. \Trayref{B}.
      \UCpaso El sistema confirma que la mascota pertenece al propietario autenticado con base en la regla \BRref{BR-PetOwnership}{Verificación de Propiedad de Mascota}. \Trayref{C}.
      \UCpaso Registra la nueva vacunación en la base de datos. \Trayref{D}.
      \UCpaso Muestra el mensaje de éxito {\bf MSG1-}``Vacuna registrada exitosamente''.
      \UCpaso Redirige al usuario a la \IUref{IU-PetDetail}{Pantalla de Detalle de Mascota}, que ahora muestra la nueva vacuna en el historial.
  \end{UCtrayectoria}
  
  %--------------------------------------
  % Trayectorias Alternativas
  %--------------------------------------
  
  \begin{UCtrayectoriaA}{A}{Datos inválidos}
      \UCpaso El sistema detecta que uno o más campos tienen un formato incorrecto (ej. fecha inválida).
      \UCpaso Muestra el mensaje de error {\bf MSG-Error} ``Por favor, corrija los campos marcados.''.
      \UCpaso[\UCactor] Oprime el botón \IUbutton{Aceptar}.
      \UCpaso[] El sistema permite al usuario corregir los datos en el formulario.
      \UCpaso Continúa en el paso 5 de la trayectoria principal.
  \end{UCtrayectoriaA}
  
  \begin{UCtrayectoriaA}{B}{Sesión inválida o expirada}
      \UCpaso El sistema detecta que el token de autenticación es inválido.
      \UCpaso Muestra el mensaje de error {\bf MSG-Error} ``No estás autorizado''.
      \UCpaso[] Redirige al usuario a la \IUref{IU-Login}{Pantalla de Inicio de Sesión}.
      \UCpaso[] Termina el caso de uso.
  \end{UCtrayectoriaA}
  
  \begin{UCtrayectoriaA}{C}{Mascota no encontrada o no autorizada}
      \UCpaso[CU05-\ref{CU05-Submit}] El sistema determina que la mascota no existe o no pertenece al usuario autenticado.
      \UCpaso Muestra el mensaje de error {\bf MSG-Error} ``Mascota no encontrada o no pertenece al propietario''.
      \UCpaso[\UCactor] Oprime el botón \IUbutton{Aceptar}.
      \UCpaso[] Termina el caso de uso.
  \end{UCtrayectoriaA}
  
  \begin{UCtrayectoriaA}{D}{Fallo inesperado del sistema}
      \UCpaso[CU05-\ref{CU05-Submit}] Ocurre un error interno en el servidor al intentar registrar la vacunación.
      \UCpaso Muestra el mensaje de error {\bf MSG-Error} ``Error al agregar la vacuna''.
      \UCpaso[\UCactor] Oprime el botón \IUbutton{Aceptar}.
      \UCpaso[] El sistema permanece en el formulario, permitiendo al usuario reintentar la operación.
  \end{UCtrayectoriaA}
  
  %--------------------------------------
  % Puntos de extensión
  %--------------------------------------
  \subsection{Puntos de extensión}
  \UCExtenssionPoint{
      % Cuando:
      El Propietario decide cancelar la operación de registro.
  }{
      % Durante la región:
      Pasos 2-4 (Mientras se encuentra en el formulario de registro).
  }{
      % Casos de uso a los que extiende:
      El Propietario oprime el botón \IUbutton{Cancelar} y el sistema regresa a \hyperlink{CU14}{CU14 Consultar Detalle de Mascota}.
  }    %-------------------------------------- TERMINA descripción del caso de uso.
